\documentclass{article}
\usepackage[T2A]{fontenc}
\usepackage[utf8]{inputenc}
\usepackage[english,russian]{babel}

% Это комментарий, LaTeX его игнорирует
% Комментарии начинаются с символа процента

\begin{document}

\section{Основы LaTeX}

Это первый абзац. LaTeX игнорирует несколько пробелов.
Один пробел,   много пробелов   — всё равно будет один пробел.

Это второй абзац. Он отделён пустой строкой.

\subsection{Специальные символы}

Чтобы напечатать специальные символы, нужно использовать экранирование:

\begin{itemize}
  \item \$ — знак доллара
  \item \% — знак процента
  \item \# — знак решётки
  \item \_ — знак подчёркивания
  \item \{ — фигурная скобка открывающая
  \item \} — фигурная скобка закрывающая
  \item \textbackslash{} — обратная косая черта
  \item \textasciitilde{} — тильда
  \item \textasciicircum{} — крышка
\end{itemize}

\subsection{Жёсткие пробелы}

Обычный пробел: номер 1.
Жёсткий пробел: номер~1 (не разорвётся при переносе строки).

\end{document}